\section{模型检验与敏感性分析}

\subsection{统一检验思路}

本文的数值结果来自同一套守恒有限体积—BDF求解框架。检验分别针对空间离散、时间积分、通量平衡和模型结构展开；这些检验用于说明计算结果在给定假设下的稳定性与内部一致性，不替代实验验证。

\subsection{空间与时间离散稳定性}

四问均进行了空间网格加密与时间积分设置检验，具体差异见各问的数值检验部分。空间加密考察径向分辨率对场值及临界时间的影响，容差和最大步长对照则考察时间积分的敏感性。应分别比较完整输出网格与题目表格取值点，避免将局部最大差异直接解释为所有表值的误差。

\subsection{守恒与结构自洽性}

共享界面通量的瞬时平衡检验有限体积离散的一致性，输出节点上的累计平衡还包含通量时间积分误差，两类指标应区分解释。各问已分别报告相应残差，本节不重复列示数值。

问题三通过与问题二前三小时结果对齐检查跨问题一致性；问题四则通过零扩散极限方程的辅助积分、固定域退化复现及代表时刻表面根扫描检查结构自洽性。其中，辅助极限积分不替代主程序完整流程检验，代表时刻扫描也不能证明全时段或全部参数下的根唯一性。

\subsection{边界与环境假设的影响}

长时段计算必须对附件1结束后的环境量作延拓。本文采用附件末值保持的基准口径；问题三的敏感性计算表明，将 $14400\,\mathrm{s}$ 后环境水分浓度提高至 $0.07\,\mathrm{kg/kg}$，临界时间约延长 $0.63\,\mathrm{h}$。这说明后期结果对外部湿度边界具有可见敏感性，临界时间应理解为给定环境延拓下的模型输出。

附件2半径轨迹的插值方法同样属于输入处理的一部分。正式模型采用PCHIP穿过全部半径节点；与分段线性插值的对照用于评估输入处理影响，不改变正文的正式结果。上述敏感性分析只讨论边界数据和插值口径，不把差异解释成实验测量误差或真实材料机理。
